\section{The chain and its elementary lower bound}\label{sec:prelim}

Let $d\ge1$ be an integer. Write $n=N=2^d$, let $V=\F^d$, and identify positions $0,\ldots,n-1$
with their binary strings, most significant bit first. Let $G=\Sym(V)$.
Fix card labels indexed by $V$, with card $x$ at position $x$ in
the identity state. A state $g\in G$ maps a label to its current
position, and permutations are composed as functions:
$(gh)(x)=g(h(x))$.

In one \emph{physical shuffle}, the deck is split into its top and bottom
halves, and, for $i=0,\ldots,n/2-1$, the cards at index $i$ within the
two halves occupy output positions
$2i$ and $2i+1$ in an independently chosen fair order. Thus one step
uses $n/2$ independent fair bits. In binary coordinates its action is
\begin{equation}\label{eq:physical-step}
 (x_1,x_2,\ldots,x_d)\longmapsto
 (x_2,\ldots,x_d,x_1+\xi_{(x_2,\ldots,x_d)}),
\end{equation}
where the bits indexed by the suffix are independent. This convention
counts individual physical shuffles, not sweeps of $d$ such shuffles.

Let $q_d$ be the probability law of the one-step position permutation,
and let $U$ be uniform measure on $G$. We also write $U_n$ or $U_G$
when indicating the state-space size or group is useful. The transition
kernel is $P_d(g,h)=q_d(hg^{-1})$; thus $P_d^t(g,\cdot)$ is the law
after $t$ shuffles started from $g$. For probability measures on a finite
set, we use
\[
 \|\mu-\nu\|_{\TV}
 =\frac12\sum_x|\mu(x)-\nu(x)|
 =\max_A|\mu(A)-\nu(A)|.
\]
For measures on $G$, the convolution $\mu*\nu$ is the law of $XY$, where
$X\sim\mu$ and $Y\sim\nu$ are independent. Starting from $g_0$, the law
after $t$ steps is the right translate of $q_d^{*t}$ by $g_0$.
Uniform measure is invariant under every such translation and under
left multiplication by every step. Consequently $U$ is stationary, and
\begin{equation}\label{eq:worst-state}
 \max_{g_0\in G}\|P_d^t(g_0,\cdot)-U\|_{\TV}
 =\|q_d^{*t}-U\|_{\TV}.
\end{equation}
Here and below $Y_t$ denotes the original shuffle started at the identity.
We define
\begin{equation}\label{eq:tmix}
 t_{\mathrm{mix}}(d)=
 \min\{t\in\mathbb Z_{\ge0}:\|q_d^{*t}-U\|_{\TV}\le1/4\}.
\end{equation}
Total variation is nonincreasing in $t$ by contraction under convolution.
The minimum exists for every $d\ge1$; the coordinate description in
Section~\ref{sec:frames} gives a short proof. Our estimates below give
the quantitative bound as $d$ grows.

\begin{lemma}[Support lower bound]\label{lem:lower}
For every integer $t\ge0$,
\begin{equation}\label{eq:exact-support-lower}
 \|q_d^{*t}-U\|_{\TV}\ge 1-\frac{2^{tn/2}}{n!}.
\end{equation}
Consequently,
\begin{equation}\label{eq:sharp-support-time}
 t_{\mathrm{mix}}(d)\ge
 \left\lceil\frac2n\log_2\left(\frac{3n!}{4}\right)\right\rceil
 =2d-O(1).
\end{equation}
For $0\le t\le d$,
\[
 \|q_d^{*t}-U\|_{\TV}\ge1-(e/\sqrt n)^n.
\]
In particular, the distance tends to one uniformly for integer $t\le d$,
and $t_{\mathrm{mix}}(d)\ge d$ for all sufficiently large $d$.
\end{lemma}
\begin{proof}
There are at most $2^{tn/2}$ choices of all shuffle bits, hence at most
that many possible states. If $A$ is their support, then
\[
 \|q_d^{*t}-U\|_{\TV}\ge q_d^{*t}(A)-U(A)
 \ge1-\frac{2^{tn/2}}{n!}.
\]
If this distance is at most $1/4$, the last ratio must be at least
$3/4$. Taking logarithms and then the least integer above the resulting
lower bound proves the first part of \eqref{eq:sharp-support-time}.
The integral bound
$\log(n!)\ge\int_1^n\log x\,dx=n\log n-n+1$
and $n!\le n^n$ give the asserted asymptotic lower bound. When
$t\le d$, the support ratio is at most $n^{n/2}/n!$, which is at most
$(e/\sqrt n)^n$. This proves the remaining statements.
\end{proof}

\begin{lemma}[One-card comparison]\label{lem:one-card-lower}
For a fixed initial card and an integer $0\le t<d$, its location law
has support at most $2^t$. Consequently the full permutation law has
total-variation distance at least $1-2^t/n$ from uniform, and
$t_{\mathrm{mix}}(d)\ge d$. After $d$ steps every individual card
has a uniform location.
\end{lemma}
\begin{proof}
At each step a card has at most two possible successors. Its support
therefore has at most $2^t$ positions. The uniform position law assigns
that support mass at most $2^t/n$, and passing from a full permutation
to the position of one card cannot increase total variation.
For the last assertion, undo the deterministic rotations. During the
first $d$ steps each coordinate is refreshed exactly once, with a bit
that is fair conditional on the entire previous card path. These
successive refreshed bits form a uniform binary string. Restoring the
rotation preserves uniformity.
\end{proof}

We shall use basic characteristic-zero representation theory of finite
groups: unitary realizations, orthogonal decomposition into irreducibles,
Schur's lemma, and matrix-coefficient orthogonality
\cite{etingof,sagan}. All representation spaces below are finite-dimensional
complex Hilbert spaces. For an operator $A$, $\|A\|_{\op}$ is its operator
norm and $\|A\|_{\HS}^2=\tr(A^*A)$ its squared Hilbert--Schmidt norm;
the trace is \emph{not} divided by the dimension. We prove explicitly the
dimension estimate and the passage from marginals to the full group law
that the argument requires.
